% !TEX TS-program = pdflatex
% !TEX encoding = UTF-8 Unicode

% This is a simple template for a LaTeX document using the "article" class.
% See "book", "report", "letter" for other types of document.

\documentclass[11pt]{article} % use larger type; default would be 10pt

\usepackage[utf8]{inputenc} % set input encoding (not needed with XeLaTeX)

%%% Examples of Article customizations
% These packages are optional, depending whether you want the features they provide.
% See the LaTeX Companion or other references for full information.

%%% PAGE DIMENSIONS
\usepackage{geometry} % to change the page dimensions
\geometry{a4paper} % or letterpaper (US) or a5paper or....
% \geometry{margin=2in} % for example, change the margins to 2 inches all round
% \geometry{landscape} % set up the page for landscape
%   read geometry.pdf for detailed page layout information

\usepackage{graphicx} % support the \includegraphics command and options

% \usepackage[parfill]{parskip} % Activate to begin paragraphs with an empty line rather than an indent

%%% PACKAGES
\usepackage{booktabs} % for much better looking tables
\usepackage{array} % for better arrays (eg matrices) in maths
\usepackage{paralist} % very flexible & customisable lists (eg. enumerate/itemize, etc.)
\usepackage{verbatim} % adds environment for commenting out blocks of text & for better verbatim
\usepackage{subfig} % make it possible to include more than one captioned figure/table in a single float
\usepackage{float}
% These packages are all incorporated in the memoir class to one degree or another...

%%% HEADERS & FOOTERS
\usepackage{fancyhdr} % This should be set AFTER setting up the page geometry
\pagestyle{fancy} % options: empty , plain , fancy
\renewcommand{\headrulewidth}{0pt} % customise the layout...
\lhead{}\chead{}\rhead{}
\lfoot{}\cfoot{\thepage}\rfoot{}

%%% SECTION TITLE APPEARANCE
\usepackage{sectsty}
\allsectionsfont{\sffamily\mdseries\upshape} % (See the fntguide.pdf for font help)
% (This matches ConTeXt defaults)

%%% ToC (table of contents) APPEARANCE
\usepackage[nottoc,notlof,notlot]{tocbibind} % Put the bibliography in the ToC
\usepackage[titles,subfigure]{tocloft} % Alter the style of the Table of Contents
\renewcommand{\cftsecfont}{\rmfamily\mdseries\upshape}
\renewcommand{\cftsecpagefont}{\rmfamily\mdseries\upshape} % No bold!
\usepackage{fancyvrb}
\usepackage{hyperref}

%%% END Article customizations

\setlength{\parindent}{0pt}

%%% The "real" document content comes below...

\title{Flight Design System}
\author{by indy91}
%\date{} % Activate to display a given date or no date (if empty),
         % otherwise the current date is printed 


\begin{document}
\maketitle

\section{Introduction}

The Flight Design System is a ground-up rendezvous mission planning tool, specifically for the Orbiter addons NASSP and SSV. Launch time and targeting parameters can be calculated and verified as well as the rendezvous profile.\\

\newpage
\section{Config}
For a new project, the first step should always be to enter a project/mission name, select the launch day and the addon (NASSP or SSV) for which the mission will be planned. NASSP does not use the standard Earth configuration file in Orbiter, so the World selection accounts for certain differences in Earth radius and rotations.\\
\newpage
\section{Generic LWP}
\newpage
\section{Shuttle LWP}
The Shuttle LWP is derived very closely from the Rendezvous Launch Window processor that the actual Flight Dynamics Officers for Space Shuttle missions used. It is described in detail in the FDO Console Handbook, see references.\\
\newpage
\section{Skylab LWP}
\newpage
\section{OMP}
\newpage
\section{State Vector}
State vectors can be added in a few formats and coordinate systems. The typical state vector for a rendezvous target will be in the TLE format, while internally stored state vectors will be in the True Equator-Greenwich (TEG) coordinate system. Also supported are J2000 and M50. While TLE data already contain information about the drag of the vehicle, the FDS uses numerical integration for trajectory propagation, so each state vector additionally needs weight, vehicle area and a K-factor for drag calculations. The K-factor is a multiplier that as used with a hardcoded drag coefficient (CD) of 2.0. In other places in the FDS the CD is entered directly.\\
\newpage
\section{Walkthrough: Shuttle rendezvous mission}
Enter project name, select SSV as world, enter Year, Month and Day of launch. Save the project. Find a TLE for the launch day, best just before liftoff time. Go to the State Vector tab, click New and enter the full TLE under SV. Change the Coordinates selection to TLE, enter weight, area and K-Factor (typically 1.0) and save. Click check to see how close to midnight of launch day the TLE time tag is.\\

Under the Shuttle LWP tab ensure all inputs are as desired. TBD\\
\newpage
\section{Reference Data}
\subsection{Radius vs. insertion altitude}

The Shuttle insertion altitude is usually referenced relative to the Earth equatorial radius of 6378166 meters. Here a convenient conversion table for radius vs. altitude.\\

\begin{tabular}{ |l|l| }
\hline
Altitude in NM & Radius in feet\\
\hline
52 & 21241700 \\
57 & 21272080 \\
60 & 21290308 \\
\hline
\end{tabular}
\newpage
\section{References}
\begin{itemize}
\item \href{https://www.ibiblio.org/apollo/Shuttle/MCC/FDO%20Console%20Handbook%20-%20On-orbit.pdf}{FDO Console Handbook, Voume 3: On-Orbit Trajectory Operations}
\end{itemize}
\end{document}
